% 仅提供必要的内容，使用方式按照自己的风格即可。
% 如果选项冲突建议在GitHub上联系作者。
% 如果在如何排版上有问题不建议联系，请认真阅读lshort-zh。
%
% 不使用宏包或者文类形式的原因：
%   1. 文类绝大部分功能都是继承，使用者修改不容易。
%   2. 宏包不易懂，同样在修改上存在障碍。
%   3. 宏包文类许多选项可以退化为命令，无需面向对象编程。
% 当然想抽成宏包或文类很简单，有兴趣可以玩一下。
%
% 控制序列若无特殊说明，提供均使用中文命令。
% 一方面防止冲突，另一方面方便理解。
% 推荐使用文类ctexart
% 如果字体加粗失败建议使用如下选项
% \documentclass[AutoFakeBold,AutoFakeSlant]{ctexart}
%
% https://github.com/Iydon
%


\makeatletter % 将@作为控制序列的一部分，可理解为私有变量
    % 使用的宏包
    \usepackage{ifxetex}
        % 建议使用XeLaTeX编译，其他编译方式建议注释此处。
        \RequireXeTeX

    \usepackage{geometry}
        % 设置页面大小
        \geometry{a4paper,top=2.54cm,bottom=2.54cm,left=3.18cm,right=3.18cm}

    \usepackage{amssymb}
        % 英文标题页的小盒子

    \usepackage{setspace}
        % 调整间距

    \usepackage{graphicx}
        % 使用图片

    \usepackage{ctex}
        % CJK排版加载宏集
        % 字号汉化设置
        \newcommand{\初号}{\zihao {0}}
        \newcommand{\小初}{\zihao{-0}}
        \newcommand{\一号}{\zihao {1}}
        \newcommand{\小一}{\zihao{-1}}
        \newcommand{\二号}{\zihao {2}}
        \newcommand{\小二}{\zihao{-2}}
        \newcommand{\三号}{\zihao {3}}
        \newcommand{\小三}{\zihao{-3}}
        \newcommand{\四号}{\zihao {4}}
        \newcommand{\小四}{\zihao{-4}}
        \newcommand{\五号}{\zihao {5}}
        \newcommand{\小五}{\zihao{-5}}
        \newcommand{\六号}{\zihao {6}}
        \newcommand{\小六}{\zihao{-6}}
        \newcommand{\七号}{\zihao {7}}
        \newcommand{\八号}{\zihao {8}}
        % 字体汉化设置
        \newcommand{\宋体}{\songti}
        \newcommand{\黑体}{\heiti}
        \newcommand{\仿宋}{\fangsong}
        \newcommand{\楷书}{\kaishu}
        \newcommand{\隶书}{\lishu}
        \newcommand{\幼圆}{\youyuan}
        \newcommand{\雅黑}{\yahei}
        \newcommand{\苹方}{\pingfang}
        % 字体效果汉化设置
        \newcommand{\粗体}{\bfseries}
        \newcommand{\斜体}{\itshape}
        % 设置章节样式
        \ctexset{
            % 需要使用ctexart文类，否则报错键不存在
            section       = {
                                name={{},.},
                                format={\hspace*{-1.5em}\黑体\三号},
                            },
            subsection    = {
                                format={\黑体\四号},
                            },
            subsubsection = {
                                format={\黑体\小四},
                            },
        }

    \usepackage{tocloft}
        % 设置目录样式
        \renewcommand\cfttoctitlefont{\hfill\黑体\小二}
        \renewcommand\cftaftertoctitle{\hfill}
        \renewcommand\cftsecfont{\粗体\宋体\四号}
        \renewcommand\cftsecpagefont{\粗体\宋体\四号}
        \renewcommand\cftsecleader{\cftdotfill{\cftdotsep}}
        \renewcommand\cftsubsecfont{\宋体\四号}
        \renewcommand\cftsubsecpagefont{\宋体\四号}
        \renewcommand\cftsubsubsecfont{\宋体\四号}
        \renewcommand\cftsubsubsecpagefont{\宋体\四号}
        \setlength\cftparskip{7pt}
        \setlength\cftsubsecindent{0pt}
        \setlength\cftsubsubsecindent{0pt}

    \usepackage{xkeyval}
        % 根据名称定义中英文控制序列
        \def\中英信息#1,#2#3{
            \expandafter\def\csname #3@中\endcsname{#1}
            \expandafter\def\csname #3@英\endcsname{#2}
        }
        % 定义键
        \newcommand{\信息中英键}[1]%
            {\define@key{信息}{#1}{\中英信息##1{#1}}}
        \信息中英键{分类号}
        \信息中英键{编号}
        \信息中英键{UDC}
        \信息中英键{密级}
        \信息中英键{题目}
        \信息中英键{子标题}
        \信息中英键{姓名}
        \信息中英键{学号}
        \信息中英键{系别}
        \信息中英键{专业}
        \信息中英键{指导老师}
        \信息中英键{时间}
        % 设置键-值
        \newcommand{\设置信息}[1]{\setkeys{信息}{#1}}


    % 目录
    \newcommand{\目录}{\setcounter{page}{1}\tableofcontents}
    % 设置标题页
    % 中文
    \newcommand\下划线[2][6em]{\hskip1pt\underline{\hb@xt@ #1{\hss#2\hss}}\hskip3pt}
    \newcommand\中文标题页{%
        \newgeometry{%
            top=2.25cm,bottom=2.25cm,
            left=2.54cm,right=2.48cm}
        \thispagestyle{empty}
        \begin{center}
            \begin{tabular}{llr}
                分类号  \下划线[2.3cm]{\分类号@中}       \hspace{0.45\textwidth} & 编号\下划线[2.3cm]{\编号@中} \\[5pt]
                U\hfill D\hfill C\下划线[2.3cm]{\UDC@中} \hspace{0.45\textwidth} & 密级\下划线[2.3cm]{\密级@中}
            \end{tabular}
            \vskip 1.35cm
            \includegraphics[width=0.85\textwidth]{logos/SUSTech_Logo_zh.pdf}
            \vskip 0.5cm
            {\粗体\宋体\小初 本科生毕业设计（论文）}
            \vskip 3.5cm
            {\粗体\宋体\三号
                \begin{tabular}{lc}
                    题\hfill 目： & \下划线[9cm]{\题目@中}     \\[9pt]
                                  & \下划线[9cm]{\子标题@中}   \\[9pt]
                    姓\hfill 名： & \下划线[9cm]{\姓名@中}     \\[9pt]
                    学\hfill 号： & \下划线[9cm]{\学号@中}     \\[9pt]
                    系\hfill 别： & \下划线[9cm]{\系别@中}     \\[9pt]
                    专\hfill 业： & \下划线[9cm]{\专业@中}     \\[9pt]
                    指导教师：    & \下划线[9cm]{\指导老师@中} \\[9pt]
                \end{tabular}}
            \vskip 3.14cm
            {\宋体\三号\时间@中}
        \end{center}
        \clearpage
        \restoregeometry
    }
    % 英文
    \newcommand\英文标题页{%
        \newgeometry{%
            top=2.54cm,bottom=2.75cm,
            left=2.54cm,right=2.48cm}
        \thispagestyle{empty}
        \begin{center}
            CLC  \下划线[2.3cm]{\分类号@英} \hfill Number\下划线[2.3cm]{\编号@英}\\[5pt]
            UDC\下划线[2.3cm]{\UDC@英} \hfill Available for reference~\quad$\square$Yes~\quad$\square$No
            \vskip 1.8cm
            \includegraphics[width=0.9\textwidth]{logos/SUSTech_Logo_en.pdf}
            \vskip 1.6cm
            {\小初 Undergraduate Thesis}
            \vskip 2.7cm
            {\粗体\三号
                \begin{tabular}{rr}
                    Thesis Title:   & \下划线[9cm]{\题目@英}     \\[9pt]
                                    & \下划线[9cm]{\子标题@英}   \\[2.1cm]
                    Student Name:   & \下划线[9cm]{\姓名@英}     \\[9pt]
                    Student ID:     & \下划线[9cm]{\学号@英}     \\[9pt]
                    Department:     & \下划线[9cm]{\系别@英}     \\[9pt]
                    Program:        & \下划线[9cm]{\专业@英}     \\[9pt]
                    Thesis Advisor: & \下划线[9cm]{\指导老师@英} \\[10pt]
                \end{tabular}}
            \vskip 2.1cm
            \三号 Date: \时间@英
        \end{center}
        \clearpage\restoregeometry
    }
    % 诚信承诺书
    % 中文
    \newcommand\中文诚信承诺书{
        \newgeometry{%
            top=3.6cm,bottom=2.75cm,
            left=2.54cm,right=2.48cm}
        \thispagestyle{empty}
        \begin{center}
            \黑体\二号 诚信承诺书
        \end{center}
        \vskip 1.8cm
        \begin{spacing}{1.9}\宋体\四号
        1.本人郑重承诺所呈交的毕业设计（论文），是在导师的指导下，独立进行研究工作所取得的成果，所有数据、图片资料均真实可靠。\par
        2.除文中已经注明引用的内容外，本论文不包含任何其他人或集体已经发表或撰写过的作品或成果。对本论文的研究作出重要贡献的个人和集体，均已在文中以明确的方式标明。\par
        3.本人承诺在毕业论文（设计）选题和研究内容过程中没有抄袭他人研究成果和伪造相关数据等行为。\par
        4.在毕业论文（设计）中对侵犯任何方面知识产权的行为，由本人承担相应的法律责任。
        \vskip 4cm
        \hfill 作者签名:\phantom{xxxx月xxxx日} \par
        \hfill\下划线[1.6cm]{}年\下划线[1cm]{}月\下划线[1cm]{}日
        \end{spacing}
        \clearpage\restoregeometry
    }
    % 英文
    \newcommand\英文诚信承诺书{
        \newgeometry{%
            top=2.75cm,bottom=2.75cm,
            left=2.54cm,right=2.48cm}
        \thispagestyle{empty}
        \begin{center}
            \huge COMMITMENT OF HONESTY
        \end{center}
        \vskip 0.9cm
        \begin{spacing}{1.45}
        \三号\noindent
        1. I solemnly promise that the paper presented comes from my independent research work under my supervisor's supervision. All statistics and images are real and reliable.\\
        2. Except for the annotated reference, the paper contents no other published work or achievement by person or group. All people making important contributions to the study of the paper have been indicated clearly in the paper.\\
        3. I promise that I did not plagiarize other people's research achievement or forge related data in the process of designing topic and research content.\\
        4. If there is violation of any intellectual property right, I will take legal responsibility myself.
        \vskip 3.63cm
        \hfill Signature:\phantom{abcdefghijklmn} \par
        \hfill Date:\phantom{abcdefghijklmnopqr}
        \end{spacing}
        \clearpage\restoregeometry
    }
    % 摘要
    \usepackage{listofitems}
    \newcommand\摘要标题{
        \begin{center}
            \黑体\二号\题目@中
            \vskip 3pt
            \hspace*{.3\textwidth}\小二 ——\子标题@中
            \vskip 1cm
            \宋体\四号\姓名@中 \\[9pt]
            （\楷书\小四\系别@中\phantom{空}指导教师： \指导老师@中）
        \end{center}
        \vskip 1.5cm
    }
    % 中文
    \newenvironment{中文摘要}[1]%
                   {\setsepchar{, }\readlist\args{#1}
                        \thispagestyle{empty}
                        \begin{spacing}{1.45}
                        \noindent\黑体\三号[摘要]：
                        \宋体\四号}%
                   {\vfill
                        \noindent\黑体\三号 [关键词]：
                        \宋体\四号 \args[1]
                        \end{spacing}
                        \vfill\clearpage}
    % 英文
    \newenvironment{英文摘要}[1]%
                   {\setsepchar{, }\readlist\args{#1}
                        \thispagestyle{empty}
                        \begin{spacing}{1.45}
                        \noindent\三号 [ABSTRACT]:
                        \四号}%
                   {\vfill
                        \noindent\三号 [Keywords]:
                        \四号 \args[1]
                        \end{spacing}
                        \vfill\clearpage}

\makeatother % 将@不作为控制序列的一部分
